\documentclass[11pt,a4paper]{article}

% ── HAL-friendly preamble ──────────────────────────────────────────
\usepackage[utf8]{inputenc}
\usepackage[T1]{fontenc}
\usepackage[english]{babel}
\usepackage[a4paper,margin=2.5cm]{geometry}
\usepackage{amsmath,amssymb,amsthm}
\usepackage{mathtools}
\usepackage{bm}
\usepackage{graphicx}
\usepackage{booktabs}
\usepackage{array}
\usepackage{caption}
\usepackage{subcaption}
\usepackage{xcolor}
\usepackage{hyperref}
\usepackage{cleveref}

\hypersetup{
  colorlinks=true,
  linkcolor=black,
  citecolor=black,
  urlcolor=blue!50!black,
  pdftitle={HCI: Harmonic Contour Integration for Edge Detection},
  pdfauthor={Author},
}

% ── Math shortcuts ──────────────────────────────────────────────────
\newcommand{\R}{\mathbb{R}}
\newcommand{\norm}[1]{\left\lVert#1\right\rVert}
\newcommand{\abs}[1]{\left\lvert#1\right\rvert}
\newcommand{\sgn}{\operatorname{sgn}}
\newcommand{\softplus}{\operatorname{softplus}}
\newcommand{\relu}{\operatorname{ReLU}}

\title{HCI: Harmonic Contour Integration for Edge Detection}

\author{%
  Siddharth Basireddy%
}

\date{\today}

\begin{document}
\maketitle

\begin{abstract}
This paper describes HCI (Harmonic Contour Integration), a compact, fully distributed edge detection method. Through the encoding of classical early-visual priors like divisive normalisation, and associative field facilitation and suppression, the method exposes a network made up of a small set of $154$ learnable parameters. The method is composed of four staged operations on an image: the construction of a pixel contrast signal (L0), harmonic projection of the signals over a pool of pixels to generate von-Mises orientation bins (L1), refinement of the signal via facilitation and suppression (Seed) and the eventual projection of the pooled signal back onto the pixel level (renderer). The renderer back projection is then passed through non-max suppression to retrieve clean edges.
\end{abstract}

\paragraph{Keywords.} Contour Integration, Harmonic Projection, Divisive Normalisation, Filtered Back-Projection, Edge Detection.

% ─────────────────────────────────────────────────────────────────────
\section{Introduction}\label{sec:intro}

One of the primary functions of the early visual cortex is the transformation of retinal input into global contours. As the cortex operates in a highly parallelized fashion, low-level signals are normalized across the visual grid and passed through layers of increasing radius that accumulate evidence across scale. These low-level signals are subjected to cortical transformations such as non-linear saturation~\cite{naka1966srod}, oriented filtering~\cite{hubel1962orientation}, divisive normalization~\cite{heeger1992normalisation,cavanaugh2002surround}, lateral facilitation and inhibition~\cite{kapadia1995dynamics,field1993contour} and figure-ground polarity~\cite{zhou2000border}. At higher levels, contour readouts are sharpened by collinear support and cross-orientation suppression before being deposited onto a coherent boundary map.

Analogs to the cortical architecture for computer vision require transforming sparse, irregular mechanisms into dense, synchronous operations suitable for array programming models and parallel processing hardware. Deep convolutional and transformer edge detectors have become the dominant approach~\cite{xie2015hed,liu2017rcf,he2019bdcn,pu2022edter,soria2023dexined,ye2024diffusionedge}, recovering contour structure implicitly through millions of learned weights that carry structural priors from the training data alone. The method introduced in this paper revisits older pipelines built from explicit structural priors, like Canny~\cite{canny1986computational} and gPb~\cite{arbelaez2011gpb}, which construct a pixel-level signal and improve it through an interpretable sequence of stages, but hold gain, scale and anisotropy fixed at design time. We deviate from these gradient-based methods, opting for harmonic projection and encoding cortex-like mechanisms as the core heuristic of a compact, fully differentiable four-stage graph, that extracts thin edges from an RGB image while exposing only $154$ learnable scalars to data.

% ─────────────────────────────────────────────────────────────────────
\section{Method}\label{sec:method}

The pipeline operates in four stages; L0, L1, seed and renderer, each with a fixed radius of interaction.

\subsection{L0: Harmonic Projection}\label{sec:l0}

Let $I:\Omega\to\R^3$ be the input image in linear RGB, range $[0,1]$. The eight-neighbour offsets at each pixel $\{\delta_k\}_{k=1}^{8}$ have bearing angles $\phi_k = (k-1)\pi/4$ in image coordinates.
We route luminance and chrominance routed through separate channels of a learned spectral metric and then squashed via a Naka-Rushton nonlinearity. This is then passed through a harmonic projection to produce per-pixel complex moments $z_1, z_2 \in \mathbb{C}$. 

\paragraph{Spectral RGB metric:}
Let $L = (1,1,1)/\sqrt{3}$ be the (fixed) achromatic axis and let $(B_1, B_2)$ be a fixed orthonormal basis of $L^{\perp}$ in $\R^3$. The metric tensor on RGB differences is
\begin{equation}\label{eq:metric}
  M = W^\top W = w_\ell^2\, L L^\top + s_1^2\, b_1 b_1^\top + s_2^2\, b_2 b_2^\top,
\end{equation}
where $(b_1, b_2)$ is the in-plane rotation of $(B_1, B_2)$ by a learned angle $\theta_{\mathrm{chr}}$:
\begin{equation}
  b_1 = \cos\theta_{\mathrm{chr}}\, B_1 + \sin\theta_{\mathrm{chr}}\, B_2,
  \qquad
  b_2 = -\sin\theta_{\mathrm{chr}}\, B_1 + \cos\theta_{\mathrm{chr}}\, B_2.
\end{equation}
The four learned scalars are $(w_\ell, s_1, s_2, \theta_{\mathrm{chr}})$ with $w_\ell, s_1, s_2 \geq 0$ via $\softplus$. Applying the rotation to the basis vectors (rather than as a left multiplier on the row-stack of $W$) is required: the latter would make $\theta_{\mathrm{chr}}$ cancel from $W^\top W$ by orthogonality and become pure gauge.

\paragraph{Directional contrasts:}
For each pixel $p$ and offset $\delta_k$, the RGB difference $\Delta c_k(p) = I(p) - I(p + \delta_k)$ is decomposed into luminance and chrominance magnitudes:
\begin{equation}\label{eq:dlumchr}
  d_k^{\mathrm{lum}}(p) = w_\ell \,\abs{L \cdot \Delta c_k(p)},
  \qquad
  d_k^{\mathrm{chr}}(p) = \sqrt{s_1^2 \bigl(b_1 \cdot \Delta c_k(p)\bigr)^2 + s_2^2 \bigl(b_2 \cdot \Delta c_k(p)\bigr)^2}.
\end{equation}

\paragraph{Naka-Rushton squash:}
Each channel is passed through a rational saturating nonlinearity~\cite{naka1966srod} with channel-specific half-saturation:
\begin{equation}\label{eq:nr}
  h_k^{\mathrm{lum}}(p) = \gamma\,
    \frac{\bigl(d_k^{\mathrm{lum}}(p)\bigr)^2}
         {\eta_{\mathrm{lum}}^{2} + \bigl(d_k^{\mathrm{lum}}(p)\bigr)^2},
  \qquad
  h_k^{\mathrm{chr}}(p) = \gamma\,
    \frac{\bigl(d_k^{\mathrm{chr}}(p)\bigr)^2}
         {\eta_{\mathrm{chr}}^{2} + \bigl(d_k^{\mathrm{chr}}(p)\bigr)^2},
\end{equation}
with $h_k(p) = h_k^{\mathrm{lum}}(p) + h_k^{\mathrm{chr}}(p)$. The constants $\eta_{\mathrm{lum}}$, $\eta_{\mathrm{chr}}$, $\gamma$ are fixed at precompute time (\cref{tab:fixed-params}); they are not gradient-bearing.

\paragraph{Optional JPEG notch:}
Before the differences in \cref{eq:dlumchr} are taken, an optional learnable separable notch filter may be applied to the projected channels $u = Wc$ to suppress the dominant JPEG block frequency. The 1D frequency response is real and even,
\begin{equation}\label{eq:notch}
  N(\omega) = 1 - d \exp\!\left(-\frac{(\abs{\omega} - \omega_n)^2}{2\sigma_n^2}\right),
\end{equation}
with three learnable scalars $(\omega_n, \sigma_n, d)$, applied as a separable cosine-synthesised kernel of half-width $H_N$ on each projected channel.

\paragraph{Harmonic projection:}
The eight values $\{h_k(p)\}_{k=1}^{8}$ are projected onto the first two even circular harmonics, yielding a complex dipole $z_1$ (polarity) and quadrupole $z_2$ (orientation, double-angle):
\begin{equation}\label{eq:zn}
  \boxed{\;
  z_n(p) = \sum_{k=1}^{8} h_k(p)\, e^{i n \phi_k},
  \qquad n \in \{1, 2\}.\;}
\end{equation}
The second-harmonic magnitude $h_{2m}(p) = \abs{z_2(p)}$ is exported separately. Image-border pixels (for which $\delta_k$ would step outside $\Omega$) are zeroed.

\subsection{L1: Cell-Level Orientation Binning}\label{sec:l1}

The image is partitioned into overlapping $P{\times}P$ cells with stride $S = P - O$, indexed by $c$. Inside cell $c$ with pixel set $\mathcal{P}_c$, the per-pixel orientation estimate is $\theta_p = \tfrac{1}{2}\arg z_2(p)$. With $K$ orientation bins of centres $\bar\theta_k = k\pi/K$ for $k=0,\dots,K-1$, the per-pixel von~Mises weight is
\begin{equation}\label{eq:vm-weight}
  g_k(\theta_p) = \exp\!\bigl(\kappa_{\mathrm{vm}}\,(\cos(2(\theta_p - \bar\theta_k)) - 1)\bigr),
\end{equation}
with learnable concentration $\kappa_{\mathrm{vm}} > 0$. The per-bin energy and sub-pixel anchors are
\begin{equation}\label{eq:rho-bin}
  \boxed{\;
  \rho^{(k)}(c) = \sum_{p \in \mathcal{P}_c} \abs{z_2(p)}\, g_k(\theta_p)\, \mathrm{ok}_p,\;}
\end{equation}
\begin{equation}\label{eq:anchor-bin}
  a_x^{(k)}(c) = \frac{\sum_{p} \abs{z_2(p)}\, g_k(\theta_p)\, x_p\, \mathrm{ok}_p}{\sum_{p} \abs{z_2(p)}\, g_k(\theta_p)\, \mathrm{ok}_p + \varepsilon},
\end{equation}
and similarly $a_y^{(k)}(c)$, where $\mathrm{ok}_p \in \{0,1\}$ excludes border pixels. Cells whose border-pixel fraction exceeds a threshold are flagged as $\mathrm{ok}(c) = 0$ and excluded from all downstream sums.

The output of L1 is the per-bin tensor $\rho^{(k)}(c) \in \R^{N \times K}$ together with the sub-pixel anchors $(a_x^{(k)}, a_y^{(k)})$ and the bin centres $\bar\theta_k$. L1 has no learned parameters of its own; $\kappa_{\mathrm{vm}}$ lives on the \textsc{Seed} module so that gradients reach it through the seed's forward pass.

\subsection{Seed: Gain Control}\label{sec:seed}

The seed refines the L1 signal, converting the per-bin energies $\rho^{(k)}$ into a per-cell scalar $\rho(c) \in [0, 1]$ and an orientation $\theta(c)$ that is sharpened by collinear facilitation along the bin tangent and suppressed by cross-orientation surround.

The per-cell state the seed produces (a boundary strength $\rho(c)$, an orientation $\theta(c)$, and a sub-pixel anchor) is deliberately task-agnostic: it is a self-contained mid-level contour representation that downstream stages could consume directly for segmentation, grouping, or shape analysis. We use the term \emph{seed} in this sense, as the distilled contour signal from which further processing proceeds. For the edge-detection task treated here that further processing is a single stage: the renderer (\cref{sec:render}) back-projects the seed onto the pixel grid. The seed itself does not depend on this choice of readout.

\paragraph{Per-bin Naka--Rushton squash:}
With $R^{(k)}(c) = \rho^{(k)}(c)\,\mathrm{ok}(c)$ and learnable half-saturation $\eta_z > 0$,
\begin{equation}\label{eq:rho-nr}
  \rho_{\mathrm{NR}}^{(k)}(c) = \frac{R^{(k)}(c)^2}{R^{(k)}(c)^2 + \eta_z^2 + \varepsilon}\,\mathrm{ok}(c).
\end{equation}

\paragraph{Same-bin collinear facilitation:}
For each bin $k$ with tangent $\hat t_k = (\cos\bar\theta_k, \sin\bar\theta_k)$ in $(dy, dx)$ order, define the directional weight
\begin{equation}\label{eq:pos-t}
  w_k(\delta) = \exp\!\left(-\frac{\norm{\delta}^2}{2\sigma_f^2}\right) \,
                \frac{(\delta \cdot \hat t_k)^2}{\norm{\delta}^2 + \varepsilon},
\end{equation}
which is large when $\delta$ is both close and aligned with the bin tangent. The collinear-pool response is
\begin{equation}\label{eq:rho-coll}
  \rho_{\mathrm{coll}}^{(k)}(c) = \relu\!\left(
  \frac{\sum_{\delta \in \mathcal{N}_f} w_k(\delta)\, \rho_{\mathrm{NR}}^{(k)}(c + \delta)}
       {\sum_{\delta \in \mathcal{N}_f} w_k(\delta) + \varepsilon}
  \right),
\end{equation}
where $\mathcal{N}_f$ is the square neighbourhood of radius $r_f$. Both $\sigma_f$ (Gaussian decay) and the bin tangent are differentiable.

\paragraph{Weighted Cross-orientation Surround:}
Let $B \in \R^{K \times K}$ be the fixed orthogonal-affinity matrix
\begin{equation}\label{eq:B}
  B_{jk} = \sin^2(\bar\theta_j - \bar\theta_k),
\end{equation}
which peaks for orthogonal bin pairs and vanishes for parallel ones. With Gaussian surround kernel $G_{\sigma_S}(\delta) = \exp(-\norm{\delta}^2 / 2\sigma_S^2)$ of radius $r_S$,
\begin{equation}\label{eq:S}
  S^{(k)}(c) = \frac{\sum_{j=1}^{K} B_{jk} \sum_{\delta \in \mathcal{N}_S} G_{\sigma_S}(\delta)\, \rho_{\mathrm{NR}}^{(j)}(c + \delta)}
                       {\bigl(\sum_{\delta} G_{\sigma_S}(\delta)\bigr) \bigl(\sum_{j} B_{jk}\bigr) + \varepsilon}.
\end{equation}
Pre-computing $\sum_{\delta} G_{\sigma_S}(\delta)\, \rho_{\mathrm{NR}}^{(j)}(c+\delta)$ as a 2D depthwise convolution and contracting with $B$ over the bin axis avoids materialising any pairwise tensor.

\paragraph{Divisive Readout:}
The per-bin excitation combines NR and collinear channels,
\begin{equation}\label{eq:e}
  e^{(k)}(c) = \beta_{\mathrm{seed}}\, \rho_{\mathrm{NR}}^{(k)}(c) + \beta_{\mathrm{coll}}\, \rho_{\mathrm{coll}}^{(k)}(c),
\end{equation}
and is divisively normalised against the surround through a Michaelis--Menten-style readout~\cite{naka1966srod,heeger1992normalisation}:
\begin{equation}\label{eq:rho-out-bin}
  \boxed{\;
  \rho_{\mathrm{out}}^{(k)}(c) = \frac{\bigl(e^{(k)}(c)\bigr)^2}
                                       {\bigl(e^{(k)}(c)\bigr)^2 + \eta_r^2 + \lambda\, \bigl(S^{(k)}(c)\bigr)^2 + \varepsilon},\;}
\end{equation}
with learnable gains $(\beta_{\mathrm{seed}}, \beta_{\mathrm{coll}})$, readout half-saturation $\eta_r$ and surround weight $\lambda$, all positive via $\softplus$. The squared $e$ on numerator and denominator gives a Naka--Rushton-shaped transfer with bounded range, while $\lambda S^2$ shifts the half-saturation point upward when the orthogonal surround is active.

\paragraph{Cell-level Export:}
Three quantities are exported per cell. The scalar boundary strength $\rho(c) \in [0, 1]$ is the maximum over bins, returned through a straight-through estimator so the forward pass is hard while the backward pass is smooth:
\begin{equation}\label{eq:rho-cell}
  \rho(c) = \max_k \rho_{\mathrm{out}}^{(k)}(c),
  \qquad
  \tilde\rho(c) = \sum_k \mathrm{softmax}_\tau\!\bigl(\rho_{\mathrm{out}}^{(\cdot)}(c)\bigr)_k\, \rho_{\mathrm{out}}^{(k)}(c),
\end{equation}
with the standard STE construction $\rho(c) - \tilde\rho(c)_{\mathrm{detach}} + \tilde\rho(c)$ on the forward graph (temperature $\tau$ fixed). The cell orientation is a double-angle weighted blend across bins,
\begin{equation}\label{eq:theta-cell}
  \theta(c) = \tfrac{1}{2}\arctan\!\left(
    \frac{\sum_k \rho_{\mathrm{out}}^{(k)}(c)\, \sin 2\bar\theta_k}
         {\sum_k \rho_{\mathrm{out}}^{(k)}(c)\, \cos 2\bar\theta_k + \varepsilon}\right),
\end{equation}
and the cell anchor is taken at the dominant bin $k^\star(c) = \arg\max_k \rho_{\mathrm{out}}^{(k)}(c)$:
\begin{equation}\label{eq:anchor-cell}
  \bigl(a_x^\star(c),\, a_y^\star(c)\bigr) = \bigl(a_x^{(k^\star)}(c),\, a_y^{(k^\star)}(c)\bigr).
\end{equation}
The seed exports the full per-bin tensor $\rho_{\mathrm{out}}^{(k)} \in \R^{N\times K}$ to the renderer, not just $\rho(c)$, so that the renderer can deposit each bin separately and aggregate with a noisy-OR.

\subsection{Renderer: Ramp-filtered Back-projection}\label{sec:render}

The renderer deposits each active (cell, bin) pair onto the pixel grid as a thin curved stroke shaped by two learned 1D kernels, then aggregates contributions via a noisy-OR. The construction borrows the geometry of filtered back-projection~\cite{kak2001ct}; each pair contributes mass along an oriented line through its anchor, and the radial filter is constrained to the windowed-ramp inverse-Radon family. This is not a filtered back-projection in the analytic sense: L0+L1 produce a local orientation-energy field rather than line integrals, the longitudinal profile $h_\parallel$ is a learned finite-support stroke rather than a ray delta, and aggregation across pairs is multiplicative (noisy-OR) rather than additive. The renderer thus exploits FBP's back-projection topology while leaving the inversion identity behind.

\paragraph{Per-bin Features and Correction MLP:}
For each (cell, bin) pair a 5-dimensional feature vector is computed from $\rho_{\mathrm{out}}^{(k)}$ and its $3{\times}3$ neighbourhood:
\begin{equation}\label{eq:feats}
  F^{(k)}_c = \bigl[\, \rho_{\mathrm{out}}^{(k)}(c),\; \mu^{(k)}_{\mathrm{coll}}(c),\; \mu^{(k)}_{t}(c),\; \textstyle\sum_{j \neq k} \rho_{\mathrm{out}}^{(j)}(c),\; \mu^{(k)}_{n}(c)\,\bigr] \in \R^5,
\end{equation}
where $\mu^{(k)}_{\mathrm{coll}}$ is a position-weighted collinear support, and $\mu^{(k)}_{t}, \mu^{(k)}_{n}$ are signed tangent and normal asymmetries of $\rho_{\mathrm{out}}^{(k)}$ across the $3{\times}3$ neighbourhood (concrete definitions in \cref{app:feats}). A small shared MLP $\Phi : \R^5 \to \R^4$ with one hidden layer of width $12$ emits four raw outputs which are bounded by tanh and softplus:
\begin{equation}\label{eq:mlp-bounded}
  \begin{aligned}
    \kappa^{(k)}_c       &= \kappa_{\max}\, \tanh\!\bigl(\Phi_1(F^{(k)}_c)\bigr), \\
    e^{(k)}_c            &= e_{\max}\, \tanh\!\bigl(\Phi_2(F^{(k)}_c)\bigr), \\
    \delta_{n}{}^{(k)}_c &= \delta_{n,\max}\, \tanh\!\bigl(\Phi_3(F^{(k)}_c)\bigr), \\
    \alpha^{(k)}_c       &= \exp\!\bigl(\alpha_{\mathrm{range}}\, \tanh\!\bigl(\Phi_4(F^{(k)}_c)\bigr)\bigr),
  \end{aligned}
\end{equation}
where $\kappa_{\max}, e_{\max}, \delta_{n,\max}, \alpha_{\mathrm{range}}$ are learnable positive bounds. The four scalars are a curvature, a tangent-direction shift of the stroke vertex, a normal-direction anchor correction, and a multiplicative amplitude $\alpha \in [e^{-\alpha_{\mathrm{range}}}, e^{+\alpha_{\mathrm{range}}}]$. At initialisation $\Phi$ outputs nearly zero, so $\kappa, e, \delta_n \approx 0$ and $\alpha \approx 1$, recovering an unmodulated straight stroke.

\paragraph{Sparsity Gate:}
Each cell's bins compete via a winner-decay gate
\begin{equation}\label{eq:gate}
  g^{(k)}(c) = \exp\!\left(\frac{\rho_{\mathrm{out}}^{(k)}(c) - \max_j \rho_{\mathrm{out}}^{(j)}(c)}{T_{\mathrm{gate}}}\right)\, \mathrm{ok}(c),
\end{equation}
with learnable temperature $T_{\mathrm{gate}} > 0$. The dominant bin always satisfies $g = 1$; lagging bins decay exponentially in the $\rho$-gap. Co-equal maxima (junctions) both fire at full gain. Pairs with $g < 10^{-3}$ or $\rho_{\mathrm{out}}^{(k)} < 10^{-4}$ are dropped from the active set, so the MLP and synthesis run only on surviving pairs.

\paragraph{Anchor correction.}
The bin anchor $(a_x^{(k)}, a_y^{(k)})$ from \cref{eq:anchor-bin} is shifted along the bin normal by $\delta_n$:
\begin{equation}\label{eq:anchor-shift}
  \tilde a^{(k)}_c = \bigl(a_x^{(k)}(c) + \delta_{n}{}^{(k)}_c \cos\bar\theta_k,\;
                            a_y^{(k)}(c) - \delta_{n}{}^{(k)}_c \sin\bar\theta_k\bigr).
\end{equation}

\paragraph{Curved Stroke Synthesis:}
For each pixel $p$ in the deposit window of $\tilde a^{(k)}_c$, let $\Delta = p - \tilde a^{(k)}_c$. Project onto the bin tangent and normal,
\begin{equation}\label{eq:proj}
  s = \Delta_y \cos\bar\theta_k + \Delta_x \sin\bar\theta_k,
  \qquad
  n = -\Delta_y \sin\bar\theta_k + \Delta_x \cos\bar\theta_k,
\end{equation}
shift the tangent coordinate by $e$ and bend the normal coordinate by the parabolic curvature term,
\begin{equation}\label{eq:bend}
  \tilde s = s - e^{(k)}_c, \qquad n_c = n - \tfrac{1}{2}\kappa^{(k)}_c\, \tilde s^{2},
\end{equation}
and synthesise the (non-negative) stroke profile by a product of learned 1D kernels:
\begin{equation}\label{eq:f}
  f^{(k)}_c(p) = \relu\!\bigl(h_{\perp}(n_c)\, h_{\parallel}(\tilde s)\bigr).
\end{equation}
Both kernels are evaluated at continuous query positions by linear interpolation between integer taps.

\paragraph{Learned 1D Kernels:}
The radial filter $h_{\perp}$ is parameterised as a windowed ramp from the inverse-Radon family,
\begin{equation}\label{eq:h-perp}
  h_{\perp}[n] = \mathrm{IFFT}\bigl(\abs{\omega}\, W(\omega)\bigr)[n],
  \qquad n \in \{-H_w, \dots, H_w\},
\end{equation}
where $W(\omega) \geq 0$ is a monotonically decreasing window in $\abs{\omega}$, parameterised by $H_w + 1$ logits $(\varphi_0, \dots, \varphi_{H_w})$ via
$W(0) = \softplus(\varphi_0)$, $W(\omega_{m+1}) = W(\omega_m)\,\sigma(\varphi_{m+1})$.
This constrains $h_{\perp}$ to the shape of an FBP ramp filter and initialises it to a Hann window of cutoff $\omega_c$; the resulting profile is the FBP-shaped piece of the renderer.

The longitudinal profile $h_{\parallel}$ is an even, monotonically decreasing function from a positive peak, parameterised by $H_w$ decay-ratio logits $(\psi_1, \dots, \psi_{H_w})$ and one peak parameter $h_0$:
\begin{equation}\label{eq:h-par}
  h_{\parallel}[m] = h_0 \prod_{j=1}^{\abs{m}} \sigma(\psi_j),
  \qquad h_0 = \softplus(h_{0,\mathrm{raw}}),
\end{equation}
initialised to a Gaussian of width $\sigma_{\parallel,0}$.

\paragraph{Per-pair Claim and Noisy-OR Aggregation:}
The per-pair pixel claim is
\begin{equation}\label{eq:claim}
  \mathrm{claim}^{(k)}_c(p) = \alpha^{(k)}_c\, g^{(k)}(c)\, \rho_{\mathrm{out}}^{(k)}(c)\, f^{(k)}_c(p),
\end{equation}
clamped into $[0, 1-\epsilon_c]$ for numerical stability, and the global boundary map is the noisy-OR over all active pairs:
\begin{equation}\label{eq:noisy-or}
  \boxed{\;
  \hat B(p) = 1 - \prod_{(c, k)\in\mathcal{A}} \bigl(1 - \mathrm{claim}^{(k)}_c(p)\bigr).\;}
\end{equation}
In log space, $\log(1 - \hat B(p)) = \sum_{(c, k)} \log(1 - \mathrm{claim}^{(k)}_c(p))$, which is computed by scatter-add on the pixel-flat index and stably back-propagates through both the synthesis and the learned kernels.

\paragraph{Inference:}
At inference, $\hat B$ is thinned by ridge non-maximum suppression along the claim-weighted dominant orientation $\theta^\star(p)$ (the circular mean of $\bar\theta_k$ weighted by per-pair claims) and thresholded at $\tau = 1/2$.

% ─────────────────────────────────────────────────────────────────────
\section{Training}\label{sec:train}

The loss is HED-style class-balanced binary cross-entropy on soft boundary maps. Let $y(p) \in [0, 1]$ be the ground-truth target and $\hat B(p) \in [0, 1]$ the prediction. With $n = \abs{\Omega}$, $n_+ = \abs{\{p : y(p) \geq \tfrac{1}{2}\}}$, $n_- = n - n_+$ and $\beta = n_-/n$, the loss is
\begin{equation}\label{eq:loss}
  \mathcal{L} = -\frac{1}{n}\sum_{p}\Bigl[\beta\, y(p)\log\hat B(p) + (1 - \beta)\,(1 - y(p))\log(1 - \hat B(p))\Bigr].
\end{equation}
Gradients flow through all four stages: through the renderer's MLP and 1D kernels, through the seed's eight scalars, through the L1 von~Mises weights via $\kappa_{\mathrm{vm}}$ (the seed runs L1 binning on the fly), and through the L0 metric scalars $(w_\ell, s_1, s_2, \theta_{\mathrm{chr}})$ and the optional notch $(\omega_n, \sigma_n, d)$. The L0 directional-difference cache stores raw $z_n$ and $h_{2m}$ tensors keyed by a version stamp; when the L0 metric is learned the cache is bypassed in favour of an on-the-fly differentiable L0 from the cached RGB image.

We train with Adam~\cite{kingma2014adam} and a cosine annealing schedule from $\eta_0 = 5\times 10^{-2}$ to $\eta_{\min} = 0.1\,\eta_0$, gradient-clipped at $\norm{\nabla}_2 \leq 1$, batch size $4$, for $15$ epochs. 

% ─────────────────────────────────────────────────────────────────────
\section{Qualitative Example}\label{sec:qualitative}

HCI applied to a single natural image is shown in \cref{fig:example}. The input (\cref{fig:base}) is a full-colour photograph with luminance and chrominance contrasts that L0 routes through the learned spectral metric. The output (\cref{fig:edges}) is the thinned boundary map obtained at inference by ridge non-maximum suppression along the claim-weighted orientation $\theta^\star(p)$ followed by thresholding at $\tau = 1/2$ (\cref{sec:render}). White pixels mark the surviving one-pixel-wide contours after thresholding the readout at 0.5.

\begin{figure}[t]
  \centering
  \begin{subfigure}[t]{0.48\textwidth}
    \vspace{0pt}
    \centering
    \includegraphics[width=\linewidth]{base.png}
    \caption{Input image $I$. The base image}
    \label{fig:base}
  \end{subfigure}%
  \hfill
  \begin{subfigure}[t]{0.48\textwidth}
    \vspace{0pt}
    \centering
    \includegraphics[width=\linewidth]{edges.png}
    \caption{Detected edges $\hat B$. Edges extracted via the pipeline and a final non-max suppression pass}
    \label{fig:edges}
  \end{subfigure}
  \caption{Qualitative result on one test image. Left: input. Right: HCI edge map at inference.}
  \label{fig:example}
\end{figure}

% ─────────────────────────────────────────────────────────────────────
\section{Parameter Inventory}\label{sec:params}

The pipeline has $154$ learnable parameters (\cref{tab:learned-params}) and a handful of fixed structural constants (\cref{tab:fixed-params}). The renderer's correction MLP accounts for $124$ of the learned parameters; everything else is interpretable scalars (gains, half-saturation constants, anisotropy widths) or small monotone-parameterised filter taps.

\begin{table}[h]
\centering
\caption{Fixed structural and numerical constants.}
\label{tab:fixed-params}
\renewcommand{\arraystretch}{1.15}
\begin{tabular}{llc}
\toprule
Block & Symbol & Value\\
\midrule
L0 NR luminance half-saturation & $\eta_{\mathrm{lum}}$ & $0.01$ \\
L0 NR chrominance half-saturation & $\eta_{\mathrm{chr}}$ & $0.01$ \\
L0 NR gain & $\gamma$ & $1.0$ \\
L0 notch half-width & $H_N$ & $4$ \\
L1 cell size & $P$ & $5$\,px \\
L1 cell overlap & $O$ & $3$\,px (stride $S = 2$) \\
L1/Seed orientation bins & $K$ & $8$ \\
Seed collinear radius & $r_f$ & $5$ cells \\
Seed surround radius & $r_S$ & $10$ cells \\
Seed STE temperature & $\tau$ & $0.1$ \\
Renderer deposit half-width & $H_w$ & $\lceil 2 S \rceil = 4$\,px \\
\bottomrule
\end{tabular}
\end{table}

\begin{table}[h]
\centering
\caption{Learned parameters. Positive scalars are reparameterised through $\softplus$; angles are unbounded.}
\label{tab:learned-params}
\renewcommand{\arraystretch}{1.15}
\begin{tabular}{llcr}
\toprule
Block & Symbol(s) & Init & \#params \\
\midrule
L0 metric gains & $w_\ell, s_1, s_2$ ($\softplus$) & $1.0$ & 3 \\
L0 chrominance angle & $\theta_{\mathrm{chr}}$ & $0$ & 1 \\
L0 notch & $\omega_n, \sigma_n, d$ & (see \cref{tab:fixed-params}) & 3 \\
\midrule
\multicolumn{3}{l}{\textbf{L0 subtotal}} & $\bm{7}$ \\
\midrule
Seed von~Mises concentration & $\kappa_{\mathrm{vm}}$ & $2.0$ & 1 \\
Seed NR half-saturation & $\eta_z$ & $10.0$ & 1 \\
Seed collinear gain & $\beta_{\mathrm{coll}}$ & $0.5$ & 1 \\
Seed channel gain & $\beta_{\mathrm{seed}}$ & $0.5$ & 1 \\
Seed readout half-saturation & $\eta_r$ & $0.3$ & 1 \\
Seed surround weight & $\lambda$ & $0.5$ & 1 \\
Seed collinear width & $\sigma_f$ & $0.5$ & 1 \\
Seed surround width & $\sigma_S$ & $2.0$ & 1 \\
\midrule
\multicolumn{3}{l}{\textbf{Seed subtotal}} & $\bm{8}$ \\
\midrule
Renderer correction MLP $5\to 12\to 4$ & $\Phi$ & near-zero output & 124 \\
Renderer correction bounds & $\kappa_{\max}, e_{\max}, \delta_{n,\max}, \alpha_{\mathrm{range}}$ & $0.1, 1, 1, 0.5$ & 4 \\
Renderer gate temperature & $T_{\mathrm{gate}}$ & $0.08$ & 1 \\
Renderer radial filter & $\{\varphi_m\}_{m=0}^{H_w}$ & Hann $\omega_c=0.5$ & 5 \\
Renderer longitudinal filter & $\{\psi_j\}_{j=1}^{H_w}, h_0$ & Gaussian $\sigma_{\parallel,0}=2$ & 5 \\
\midrule
\multicolumn{3}{l}{\textbf{Renderer subtotal}} & $\bm{139}$ \\
\midrule
\multicolumn{3}{l}{\textbf{Total}} & $\bm{154}$ \\
\bottomrule
\end{tabular}
\end{table}

% ─────────────────────────────────────────────────────────────────────
\section{Discussion}\label{sec:discussion}

HCI compresses a classical contour-integration computation into a fully differentiable graph of $154$ scalars. The choice of what to expose to data and what to fix, is rooted in the strict interpretation of cortical mechanisms. The renderer's $5 \to 12 \to 4$ correction head is the smallest learned component that empirically suffices to handle off-anchor edges, sub-pixel curvature, and bin selectivity at junctions; replacing it with a fixed identity recovers a straight-stroke depositor that is competitive but coarser. The kernels $h_{\perp}, h_{\parallel}$ are parameterised so that every point in parameter space is a valid windowed-ramp radial filter and a positive monotone-decay longitudinal profile, respectively. This keeps the synthesis path interpretable as a ramp-filtered stroke deposit, while leaving the kernel shape open to gradient descent against the boundary objective.

This report fixes the equations; experimental evaluation, ablations and comparison against deep baselines~\cite{xie2015hed,liu2017rcf,he2019bdcn,pu2022edter,soria2023dexined,ye2024diffusionedge} are deferred to a forthcoming full paper.

\section*{Code availability}
The reference implementation is available at \url{https://github.com/RenderBear/hci}.

% =====================================================================
\appendix

\section{Per-bin Renderer Features}\label{app:feats}

The five features in \cref{eq:feats} are computed from the 3$\times$3 neighbourhood $\mathcal{N}_3(c)$ of cell $c$ (with $\delta \in \mathcal{N}_3 \setminus \{0\}$) as follows. Let $t_k(\delta) = \delta_y \cos\bar\theta_k + \delta_x \sin\bar\theta_k$ and $n_k(\delta) = -\delta_y \sin\bar\theta_k + \delta_x \cos\bar\theta_k$ be the tangent and normal projections of the offset in bin $k$, and $\mathrm{pos}_t^{(k)}(\delta) = t_k(\delta)^2/\norm{\delta}^2$. Then
\begin{align}
  \mu^{(k)}_{\mathrm{coll}}(c) &= \frac{\sum_\delta \rho_{\mathrm{out}}^{(k)}(c+\delta)\,\mathrm{pos}_t^{(k)}(\delta)}
                                       {\sum_\delta \mathrm{pos}_t^{(k)}(\delta) + \varepsilon'}, \\
  \mu^{(k)}_t(c) &= \frac{\sum_\delta \rho_{\mathrm{out}}^{(k)}(c+\delta)\, t_k(\delta)}
                          {\sum_\delta \rho_{\mathrm{out}}^{(k)}(c+\delta)\, \abs{t_k(\delta)} + \varepsilon'}, \\
  \mu^{(k)}_n(c) &= \frac{\sum_\delta \rho_{\mathrm{out}}^{(k)}(c+\delta)\, n_k(\delta)}
                          {\sum_\delta \rho_{\mathrm{out}}^{(k)}(c+\delta)\, \abs{n_k(\delta)} + \varepsilon'},
\end{align}
with a soft floor $\varepsilon' = 5\times 10^{-2}$ on the denominator. $\mu^{(k)}_{\mathrm{coll}}$ summarises collinear support along the bin tangent; $\mu^{(k)}_t \in [-1, 1]$ is a signed tangent asymmetry that signals where, along the tangent, the local mass is concentrated and drives $e$; $\mu^{(k)}_n \in [-1, 1]$ is the analogous normal asymmetry and drives $\delta_n$.

\section{Renderer parameter count}\label{app:render-count}

The correction MLP has $5 \cdot 12 + 12 = 72$ parameters in the input layer and $12 \cdot 4 + 4 = 52$ in the output layer, for $124$ total. The kernel radial logits $\{\varphi_m\}$ have $H_w + 1 = 5$ parameters; the longitudinal logits $\{\psi_j\}$ have $H_w = 4$ parameters; the longitudinal peak adds $1$. The four correction bounds $(\kappa_{\max}, e_{\max}, \delta_{n,\max}, \alpha_{\mathrm{range}})$ and the gate temperature $T_{\mathrm{gate}}$ contribute $5$ further scalars. Total: $124 + 4 + 1 + 5 + 4 + 1 = 139$.

% =====================================================================
\begin{thebibliography}{99}

\bibitem{naka1966srod}
K.~I.~Naka, W.~A.~H.~Rushton.
\newblock S-potentials from colour units in the retina of fish (Cyprinidae).
\newblock \emph{Journal of Physiology}, 185(3):536--555, 1966.

\bibitem{hubel1962orientation}
D.~H.~Hubel, T.~N.~Wiesel.
\newblock Receptive fields, binocular interaction and functional architecture in the cat's visual cortex.
\newblock \emph{Journal of Physiology}, 160(1):106--154, 1962.

\bibitem{heeger1992normalisation}
D.~J.~Heeger.
\newblock Normalization of cell responses in cat striate cortex.
\newblock \emph{Visual Neuroscience}, 9(2):181--197, 1992.

\bibitem{cavanaugh2002surround}
J.~R.~Cavanaugh, W.~Bair, J.~A.~Movshon.
\newblock Nature and interaction of signals from the receptive field center and surround in macaque V1 neurons.
\newblock \emph{Journal of Neurophysiology}, 88(5):2530--2546, 2002.

\bibitem{kapadia1995dynamics}
M.~K.~Kapadia, M.~Ito, C.~D.~Gilbert, G.~Westheimer.
\newblock Improvement in visual sensitivity by changes in local context: parallel studies in human observers and in V1 of alert monkeys.
\newblock \emph{Neuron}, 15(4):843--856, 1995.

\bibitem{field1993contour}
D.~J.~Field, A.~Hayes, R.~F.~Hess.
\newblock Contour integration by the human visual system: Evidence for a local ``association field''.
\newblock \emph{Vision Research}, 33(2):173--193, 1993.

\bibitem{zhou2000border}
H.~Zhou, H.~S.~Friedman, R.~von~der~Heydt.
\newblock Coding of border ownership in monkey visual cortex.
\newblock \emph{Journal of Neuroscience}, 20(17):6594--6611, 2000.

\bibitem{canny1986computational}
J.~Canny.
\newblock A computational approach to edge detection.
\newblock \emph{IEEE Transactions on Pattern Analysis and Machine Intelligence}, 8(6):679--698, 1986.

\bibitem{arbelaez2011gpb}
P.~Arbeláez, M.~Maire, C.~Fowlkes, J.~Malik.
\newblock Contour detection and hierarchical image segmentation.
\newblock \emph{IEEE Transactions on Pattern Analysis and Machine Intelligence}, 33(5):898--916, 2011.

\bibitem{kak2001ct}
A.~C.~Kak, M.~Slaney.
\newblock \emph{Principles of Computerized Tomographic Imaging}.
\newblock SIAM Classics in Applied Mathematics 33, 2001.

\bibitem{kingma2014adam}
D.~P.~Kingma, J.~Ba.
\newblock Adam: A method for stochastic optimization.
\newblock In \emph{International Conference on Learning Representations (ICLR)}, 2015.

\bibitem{xie2015hed}
S.~Xie, Z.~Tu.
\newblock Holistically-Nested Edge Detection.
\newblock In \emph{IEEE International Conference on Computer Vision (ICCV)}, 2015.

\bibitem{liu2017rcf}
Y.~Liu, M.-M.~Cheng, X.~Hu, K.~Wang, X.~Bai.
\newblock Richer Convolutional Features for Edge Detection.
\newblock In \emph{IEEE Conference on Computer Vision and Pattern Recognition (CVPR)}, 2017.

\bibitem{he2019bdcn}
J.~He, S.~Zhang, M.~Yang, Y.~Shan, T.~Huang.
\newblock BDCN: Bi-Directional Cascade Network for Perceptual Edge Detection.
\newblock In \emph{IEEE Conference on Computer Vision and Pattern Recognition (CVPR)}, 2019.

\bibitem{pu2022edter}
M.~Pu, Y.~Huang, Y.~Liu, Q.~Guan, H.~Ling.
\newblock EDTER: Edge Detection with Transformer.
\newblock In \emph{IEEE Conference on Computer Vision and Pattern Recognition (CVPR)}, 2022.

\bibitem{soria2023dexined}
X.~Soria, A.~Sappa, P.~Humanante, A.~Akbarinia.
\newblock Dense Extreme Inception Network for Edge Detection.
\newblock \emph{Pattern Recognition}, 139:109461, 2023.

\bibitem{ye2024diffusionedge}
Y.~Ye, K.~Xu, Y.~Huang, R.~Yi, Z.~Cai.
\newblock DiffusionEdge: Diffusion Probabilistic Model for Crisp Edge Detection.
\newblock In \emph{AAAI Conference on Artificial Intelligence}, 2024.

\end{thebibliography}

\end{document}